% docs/texte/10_engineering/design.tex -- chapter "Basic Design" (Interferosim)
% Version: 0.1.2 (2026-10-12) -- first draft of the software structure (UML);
%   0.1.2: note on the implemented package sky;
%   0.1.1: unnumbered paragraphs, pipeline diagram on a landscape page
% ================================================================

\index{basic design}\index{UML}
This chapter describes the first basic design of the package
\texttt{interferosim}. It is a \emph{draft}: no simulation code has been
migrated yet. The classes are derived from the existing sources --
\texttt{aas6sim}~1.1.0 (\texttt{model.Component}, \texttt{model.Instrument},
\texttt{visibility.VisData}) and the teaching notebooks A (Wiener filter
and CLEAN) and B (the Sun as a first target) -- and will be refined in
phase~1. All diagrams use UML~2.5 notation; attributes and operations are
given with Python types, \underline{underlined} operations are class
methods, \textit{italic} operations are abstract.

\section{Package structure}

Figure~\ref{fig:uml-packages} shows the packages and their dependencies.
The design follows a layered structure:

\begin{itemize}
  \item \textbf{Core} (\texttt{sky}, \texttt{instrument},
        \texttt{observation}, \texttt{imaging}, \texttt{deconvolution}):
        the complete processing chain. The core does not depend on any
        user-interface package, so a later desktop front end can use it
        unchanged.
  \item \textbf{User interface} (\texttt{plotting}, \texttt{widgets}):
        uses the core, never the other way round. The notebooks only call
        these packages and the core.
  \item \textbf{Support} (\texttt{io}, \texttt{constants}): FITS
        input/output taken over from \texttt{aas6sim}, physical constants
        and line catalogue.
  \item \textbf{Configuration} (\texttt{configs/*.toml}): antenna types,
        receivers and sites are data, not code. A new antenna type that
        fits an existing class needs only a new configuration file.
\end{itemize}

The dependencies inside the core follow the signal path: sky and
instrument $\rightarrow$ observation $\rightarrow$ imaging
$\rightarrow$ deconvolution. There are no cyclic dependencies.

\begin{figure}[htbp]
  \centering
  % docs/texte/10_engineering/uml/packages.tex -- UML package diagram (basic design)
% Version: 0.1.0 (2026-10-12) -- first draft, based on the approved repository structure
\begin{tikzpicture}[node distance=1cm]
  % ---- user interface layer ------------------------------------------
  \umlpkg{nb}{0,0}{notebooks}
  \umlpkg{widgets}{-2.6,-2.3}{widgets}
  \umlpkg{plotting}{2.6,-2.3}{plotting}

  % ---- simulation core (independent of the user interface) -----------
  \umlpkg{sky}{-4.6,-4.9}{sky}
  \umlpkg{instrument}{-1.5,-4.9}{instrument}
  \umlpkg{observation}{-3.05,-7.1}{observation}
  \umlpkg{imaging}{0.25,-7.1}{imaging}
  \umlpkg{deconvolution}{3.55,-7.1}{deconvolution}
  \umlframe[fill=gray!8]{core}{(sky)(instrument)(observation)(imaging)(deconvolution)}{core (UI-independent)}

  % ---- support -------------------------------------------------------
  \umlpkg{io}{-3.05,-10.0}{io}
  \umlpkg{constants}{0.25,-10.0}{constants}
  \umlframe{pkg}{(widgets)(plotting)(core)(io)(constants)}{interferosim}

  % ---- configuration files (artefact, outside the package) -----------
  \node[umlpkgbox,fill=yellow!12,font=\scriptsize\sffamily,minimum width=2.2cm] (cfg)
        at (6.6,-4.9) {{\tiny$\ll$artifact$\gg$}\\\textbf{configs/*.toml}};

  % ---- dependencies --------------------------------------------------
  \draw[depend] (nb) -- node[stereo,left=2pt]{$\ll$use$\gg$} (widgets);
  \draw[depend] (nb) -- (plotting);
  \draw[depend] (widgets) -- node[stereo,above]{$\ll$use$\gg$} (plotting);
  \draw[depend] (widgets.south) -- (widgets.south |- core.north)
        node[stereo,pos=0.45,left=2pt]{$\ll$use$\gg$};
  \draw[depend] (plotting.south) -- (plotting.south |- core.north)
        node[stereo,pos=0.45,right=2pt]{$\ll$use$\gg$};

  \draw[depend] (observation) -- (sky);
  \draw[depend] (observation) -- (instrument);
  \draw[depend] (imaging) -- (observation);
  \draw[depend] (deconvolution) -- (imaging);
  \draw[depend] (instrument) -- node[stereo,above]{$\ll$read$\gg$} (cfg);
  \draw[depend] (observation) -- node[stereo,left]{$\ll$create$\gg$ VisData} (io);
  \draw[depend] (imaging) -- node[stereo,right=1pt]{$\ll$read$\gg$} (io);
  \draw[depend] (core.south -| constants) -- node[stereo,right=2pt]{$\ll$use$\gg$} (constants);
\end{tikzpicture}

  \caption{UML package diagram of \texttt{interferosim} (draft 0.1)}
  \label{fig:uml-packages}
\end{figure}

\section{Sky and instrument}

Figure~\ref{fig:uml-model} shows the model classes.

\paragraph*{Sky.} \texttt{aas6sim} describes every source with one class
\texttt{Component} and a \texttt{kind} field (\texttt{"point"},
\texttt{"gauss"}). The design replaces this by an abstract base class
\texttt{SkyComponent} with one subclass per source type. Every subclass
provides its own analytic \texttt{visibility(u, v)}; this is needed for
the solar disc from notebook~B, whose visibility is a Hankel transform
of the brightness profile (uniform or limb-brightened). A
\texttt{SkyModel} is composed of one or more components; \texttt{demo()}
returns the demo field of \texttt{aas6sim}. \emph{Implemented in version 0.2.0}
(chapter~\ref{ch:sky-models}); the positions are called
\texttt{dx\_deg}/\texttt{dy\_deg} as in the AAS6 file format, and the
field centre is a separate class \texttt{FieldCenter}.

\paragraph*{Instrument.} \texttt{aas6sim.Instrument} combines antenna,
receiver, site and baseline in a single class. The design separates
them, because several antenna types must be selectable:
\begin{itemize}
  \item \texttt{Antenna} (abstract) with the subclasses \texttt{Horn}
        and \texttt{ParabolicDish}; further types are added as
        subclasses.
  \item \texttt{Receiver} with the device profile of the Ettus USRP B210
        (\texttt{configs/receivers/b210.toml}).
  \item \texttt{Site}, \texttt{Station} (antenna at a position) and
        \texttt{Array} (site, stations, receiver), which provides the
        baselines.
\end{itemize}

\begin{figure}[htbp]
  \centering
  \resizebox{\textwidth}{!}{% docs/texte/10_engineering/uml/classes_model.tex -- UML class diagram: sky and instrument (basic design)
% Version: 0.4.0 (2026-10-12) -- package instrument as implemented (interferosim 0.3.0)
%   0.3.0: spectrum attribute, Yagi, Station with geodetic position
%   0.2.0: package sky as implemented (interferosim 0.2.0)
% Origin of the classes: aas6sim 1.1.0 (model.Component, model.Instrument)
% and notebook B (disc profiles, horn); see chapter "Basic Design".
\begin{tikzpicture}
  % =================== package sky ===================================
  % 0.2.0: as implemented in interferosim 0.2.0 (JSON format of AAS6
  % sim_sky.json: dx_deg/dy_deg, field centre, kind "disc" added)
  \umlabstract{comp}{0,0}{SkyComponent}
    {+ name : str\\+ dx\_deg : float\\+ dy\_deg : float\\+ flux\_jy : float\\+ spectrum : Spectrum\\+ \underline{kind : str}}
    {+ \textit{render(x, y, pix\_deg, out)}\\+ \textit{visibility(u, v) : ndarray}\\+ to\_dict() : dict}
  \umlclass{point}{-3.4,-3.6}{PointSource}{}{}
  \umlclass{gauss}{0,-3.6}{GaussianSource}
    {+ fwhm\_x\_deg : float\\+ fwhm\_y\_deg : float\\+ pa\_deg : float}{}
  \umlclass{disc}{3.6,-3.6}{DiscSource}
    {+ diameter\_deg : float\\+ limb : float\\+ limb\_power : float}
    {+ profile(rho) : ndarray}
  \umlclass{skymodel}{-4.9,0}{SkyModel}
    {+ name : str\\+ note : str\\+ center : FieldCenter\\+ components : list}
    {+ add(c : SkyComponent)\\+ render\_map(fov, pix) : ndarray\\+ visibility(u, v) : ndarray\\+ to\_json(path)\\+ \underline{from\_json(path) : SkyModel}\\+ \underline{demo() : SkyModel}}
  \umldata{center}{-8.9,0}{FieldCenter}
    {+ label : str\\+ ra\_deg : float\\+ dec\_deg : float}{}

  \draw[inherit] (point.north) |- ++(0,0.55) -| (comp.south);
  \draw[inherit] (gauss.north) -- (comp.south);
  \draw[inherit] (disc.north) |- ++(0,0.55) -| (comp.south);
  \draw[compose] (skymodel.east) -- node[mult,above,pos=0.8]{1..*} (comp.west);
  \draw[compose] (skymodel.west) -- node[mult,above,pos=0.75]{1} (center.east);

  \umlframe{skypkg}{(comp)(point)(gauss)(disc)(skymodel)(center)}{sky}

  % =================== package instrument ============================
  \umlabstract{ant}{0,-10.3}{Antenna}
    {+ name : str}
    {+ \textit{gain\_dbi(f\_hz) : float}\\+ effective\_area(f\_hz) : float\\+ hpbw\_deg(f\_hz) : float\\+ beam(theta\_deg, f\_hz) : ndarray\\+ jy\_to\_k(f\_hz) : float\\+ \underline{from\_toml(path) : Antenna}}
  \umlclass{horn}{-2.0,-14.4}{Horn}{+ gain\_dbi\_nominal : float\\+ design\_freq\_hz : float}{}
  \umlclass{yagi}{6.4,-14.4}{Yagi}
    {+ n\_elements : int\\+ boom\_length\_m : float\\+ gain\_dbi\_nominal : float\\+ design\_freq\_hz : float}{}
  \umlclass{dish}{2.2,-14.4}{ParabolicDish}
    {+ diameter\_m : float\\+ aperture\_efficiency : float\\+ beam\_factor : float}
    {+ u\_max(f\_hz) : float\\+ u\_cut\_eff(f\_hz) : float\\+ otf(u, f\_hz) : ndarray}
  \umlclass{station}{-4.2,-11.6}{Station}
    {+ name : str\\+ lat\_deg, lon\_deg : float\\+ height\_m : float\\+ antenna : Antenna}
    {+ ecef\_m() : ndarray\\+ enu\_m(site) : ndarray\\+ \underline{from\_enu(...) : Station}}
  \umlclass{array}{-4.0,-7.7}{Array}
    {+ site : Site\\+ stations : list\\+ receiver : Receiver}
    {+ baselines() : ndarray\\+ baselines\_enu() : ndarray\\+ baseline\_names() : list\\+ \underline{from\_toml(path) : Array}}
  \umlclass{site}{-7.9,-7.7}{Site}
    {+ name : str\\+ lat\_deg : float\\+ lon\_deg : float\\+ alt\_m : float}{+ ecef\_m() : ndarray\\+ enu\_to\_ecef(e, n, u) : ndarray}
  \umlclass{recv}{-8.6,-12.0}{Receiver}
    {+ name : str\\+ t\_sys\_k : float\\+ bandwidth\_hz : float\\+ sample\_rate\_sps : float\\+ rx\_channels : int}
    {+ sefd\_jy(a : Antenna, f\_hz) : float\\+ baseline\_sigma\_jy(a1, a2, t, f\_hz) : float\\+ \underline{from\_toml(path) : Receiver}}

  \draw[inherit] (horn.north) |- ++(0,0.5) -| (ant.south);
  \draw[inherit] (dish.north) |- ++(0,0.5) -| (ant.south);
  \draw[inherit] (yagi.north) |- ++(0,0.5) -| (ant.south);
  \draw[compose] (array.south) -- node[mult,right,pos=0.75]{1..*} (station.north);
  \draw[assoc] (station.east) -- node[mult,above,pos=0.75]{1} (ant.west);
  \draw[assoc] (array.west) -- node[mult,above,pos=0.75]{1} (site.east);
  \draw[assoc] (array.south west) -- node[mult,above,pos=0.8]{1} (recv.north east);

  \umlframe{instpkg}{(ant)(horn)(dish)(yagi)(station)(array)(site)(recv)}{instrument}
\end{tikzpicture}
}
  \caption{UML class diagram: packages \texttt{sky} (as implemented in
           version 0.2.0) and \texttt{instrument} (draft 0.1)}
  \label{fig:uml-model}
\end{figure}

\section{Observation, imaging and deconvolution}

Figure~\ref{fig:uml-pipeline} shows the processing chain.

\paragraph*{Observation.} \texttt{Observation} holds the observing
parameters (array, frequency, declination, hour-angle range,
integration time, elevation limit). \texttt{UVCoverage} is derived from
it. \texttt{VisibilitySimulator} combines sky model and observation and
creates a \texttt{VisData} object -- the data class of
\texttt{aas6sim.visibility}, so simulated and later measured data share
one format. \texttt{Correlator} covers the adding and the multiplying
interferometer, fringe rate and delay (notebook~B).

\paragraph*{Imaging.} \texttt{Imager} grids the visibilities with a
\texttt{Gridder} and applies an exchangeable \texttt{Weighting}
(strategy pattern): natural, uniform or Wiener weighting.

\paragraph*{Deconvolution.} The interface \texttt{Deconvolver} is
realised by \texttt{HogbomClean} and \texttt{WienerFilter}; both return
a \texttt{DeconvResult} (model, residual, restored image).
\texttt{CleanBeam} is fitted to the main lobe of the dirty beam.

The Wiener filter appears twice on purpose. Notebook~A shows that for a
sparsely sampled interferometer it only re-weights the measured uv
cells -- it is therefore modelled as \texttt{WienerWeighting}. For a
filled aperture (the 3\,m dish) it is a real deconvolution and is
modelled as \texttt{WienerFilter}.

% landscape page: the pipeline diagram is too wide for portrait
\begin{sidewaysfigure}
  \centering
  \resizebox{!}{14cm}{% docs/texte/10_engineering/uml/classes_pipeline.tex -- UML class diagram: observation, imaging, deconvolution (basic design)
% Version: 0.3.0 (2026-10-31) -- Observation and VisData as implemented (interferosim 0.5.0)
%   0.2.0 (2026-10-12): VisData in module visdata, frequency axis
%   0.1.0: first draft
% Origin: aas6sim 1.1.0 (visibility.VisData, model.uv_track_ew, model.elevation_deg)
% and notebooks A/B (gridding, weightings, hogbom_fft, Wiener filter, fringe rate).
\begin{tikzpicture}
  % ----- referenced classes from other packages (name only) ----------
  \node[umlpkgbox,minimum width=2cm,dashed] (rsky) at (-6.2,1.0) {sky::SkyModel};
  \node[umlpkgbox,minimum width=2cm,dashed] (rarr) at (-6.2,-1.2) {instrument::Array};

  % =================== package observation ==========================
  \umldata{obs}{-2.6,-1.2}{Observation}
    {+ array : Array\\+ ra\_deg, dec\_deg : float\\+ freq\_hz : float\\+ ha\_range\_h : tuple\\+ t\_int\_s : float\\+ min\_elev\_deg : float\\+ date : str}
    {+ hour\_angles() : ndarray\\+ elevation\_deg(ha) : ndarray\\+ uvw(ha) : ndarray\\+ delay\_s(ha) : ndarray\\+ fringe\_rate\_hz(ha) : ndarray}
  \umldata{uvc}{1.9,-1.2}{UVCoverage}
    {+ u, v, w : ndarray\\+ ha\_deg : ndarray}
    {+ \underline{from\_observation(o) : UVCoverage}}
  \umlclass{sim}{-2.6,3.2}{VisibilitySimulator}
    {+ add\_noise : bool\\+ seed : int}
    {+ run(sky, obs) : VisData}
  \umlclass{corr}{6.8,-1.2}{Correlator}
    {+ mode : \{add, multiply\}}
    {+ fringe\_rate\_hz(obs, ha) : ndarray\\+ delay\_ns(obs, ha) : ndarray\\+ correlate(x1, x2) : ndarray}
  \umlframe{obspkg}{(obs)(uvc)(sim)(corr)}{observation}

  \umldata{vis}{11.6,-1.2}{VisData}
    {+ time\_mjd, ha\_deg : ndarray\\+ u, v, w : ndarray\\+ baseline : ndarray\\+ freq\_hz : ndarray\\+ vis : ndarray (n\_rows, n\_ch)\\+ weight : ndarray\\+ header : dict}{}
  \umlframe[dashed]{iopkg}{(vis)}{visdata}

  \draw[aggregate] (obs.west) -- node[mult,above,pos=0.8]{1} (rarr.east);
  \draw[depend] (sim.west) -- node[stereo,above]{$\ll$use$\gg$} (rsky.east);
  \draw[depend] (sim.south) -- node[stereo,right]{$\ll$use$\gg$} (obs.north);
  \draw[depend] (sim.east) -| node[stereo,above,pos=0.25]{$\ll$create$\gg$} (vis.north);
  \draw[depend] (uvc.west) -- node[stereo,above]{$\ll$use$\gg$} (obs.east);

  % =================== package imaging ==============================
  \umlinterface{wgt}{-3.2,-6.4}{Weighting}{}{+ weights(W : ndarray) : ndarray}
  \umlclass{nat}{-6.3,-8.8}{NaturalWeighting}{}{}
  \umlclass{uni}{-3.2,-8.8}{UniformWeighting}{}{}
  \umlclass{wien}{-0.1,-8.8}{WienerWeighting}{+ k : float}{}
  \umlclass{grid}{-7.4,-6.4}{Gridder}
    {+ npix : int\\+ cell\_arcsec : float}
    {+ grid(v : VisData) : (V, W)}
  \umlclass{imager}{1.6,-6.4}{Imager}
    {+ gridder : Gridder\\+ weighting : Weighting}
    {+ dirty\_image(v : VisData) : Image\\+ dirty\_beam(v : VisData) : Image}
  \umldata{img}{6.0,-6.4}{Image}
    {+ data : ndarray\\+ pix\_arcsec : float\\+ unit : str}{}

  \draw[realize] (nat.north) |- ++(0,0.5) -| (wgt.south);
  \draw[realize] (uni.north) -- (wgt.south);
  \draw[realize] (wien.north) |- ++(0,0.5) -| (wgt.south);
  \draw[compose] (imager.north west) to[out=150,in=30] node[mult,above,pos=0.85]{1} (grid.north east);
  \draw[aggregate] (imager.west) -- node[mult,above,pos=0.75]{1} (wgt.east);
  \draw[depend] (imager.east) -- node[stereo,above]{$\ll$create$\gg$} (img.west);
  \umlframe{imgpkg}{(wgt)(nat)(uni)(wien)(grid)(imager)(img)}{imaging}

  \draw[depend] (imager.north) -- ++(0,1.15) -| node[stereo,above,pos=0.25]{$\ll$read$\gg$} (vis.south);

  % =================== package deconvolution ========================
  \umlinterface{dec}{1.6,-12.2}{Deconvolver}{}{+ run(dirty : Image, beam : Image) : DeconvResult}
  \umlclass{hog}{-1.6,-15.0}{HogbomClean}
    {+ gain : float\\+ n\_iter : int\\+ threshold : float}{}
  \umlclass{wf}{3.4,-15.0}{WienerFilter}{+ nsr : float}{}
  \umlclass{cb}{-6.0,-15.0}{CleanBeam}
    {+ fwhm\_arcsec : float}
    {+ \underline{fit(beam : Image) : CleanBeam}\\+ convolve(img : Image) : Image}
  \umldata{res}{6.6,-12.2}{DeconvResult}
    {+ model : Image\\+ residual : Image\\+ restored : Image}{}

  \draw[realize] (hog.north) |- ++(0,0.55) -| (dec.south);
  \draw[realize] (wf.north) |- ++(0,0.55) -| (dec.south);
  \draw[depend] (hog.west) -- node[stereo,above]{$\ll$use$\gg$} (cb.east);
  \draw[depend] (dec.east) -- node[stereo,above]{$\ll$create$\gg$} (res.west);
  \umlframe{decpkg}{(dec)(hog)(wf)(cb)(res)}{deconvolution}

  \draw[depend] (dec.north) -- node[stereo,right]{$\ll$use$\gg$} (dec.north |- imgpkg.south);
\end{tikzpicture}
}
  \caption{UML class diagram: packages \texttt{observation},
           \texttt{imaging}, \texttt{deconvolution} and the data class
           \texttt{io.VisData} (draft 0.1)}
  \label{fig:uml-pipeline}
\end{sidewaysfigure}

\section{Open points}

The following points are not yet decided and will be settled in
phase~1:
\begin{itemize}
  \item Spectral dimension: HI and OH line sources need several
        frequency channels; the draft models a single frequency.
  \item Time-domain simulation of a transit (notebook~B): own class or
        part of \texttt{VisibilitySimulator}.
  \item Further antenna types besides horn and parabolic dish.
  \item Whether \texttt{VisData} stays in \texttt{io} (as in
        \texttt{aas6sim}) or moves to \texttt{observation}.
\end{itemize}
